\section{A PTAS for uniform shallow caps in convex position}
\label{sec:uniform-cap-ptas}

The bounded-cap restriction also permits an approximation scheme for
uniform expected depth when the hypotheses are in convex position.
The construction first clears short gaps in their cyclic order.
Every remaining query then acts inside one bounded-size block.
A count-vector dynamic program retains all choices of how to interleave
queries from these blocks. Uniform weights make blocks with the same
local test family interchangeable.

\begin{theorem}\label{thm:uniform-cap-ptas}
Suppose $\HH=P\cup\{p\}$ is in strictly convex position, $p$ is a
distinguished hypothesis, and every supplied side missing $p$ contains
at most a fixed number $b\ge1$ of hypotheses. Under the uniform prior,
expected depth admits a polynomial-time approximation scheme. Every
query in the returned tree is supplied, and the result permits input
families with no singleton caps.

For $0<\varepsilon\le1$, the returned cost is at most
$(1+\varepsilon)\OPT_E$. For the block length $\ell$ chosen in the
final calibration, put
$d_{b,\ell}=2^{(b+1)\ell}$. One running-time bound is
$f(b,\ell)(n+1)^{d_{b,\ell}}\operatorname{poly}(n,m,L)$, where
$L$ is the input bit length and $f$ depends only on $b,\ell$.
The polynomial exponent depends on $b$ and $\varepsilon$; the
guarantee is a PTAS.
\end{theorem}

\begin{proofstrategy}
\begin{enumerate}
\item Clear short gaps in the convex cyclic order using supplied
caps containing marked points. A cap cannot reach two remaining blocks.
\item Describe each small block by its supplied interval masks.
Equal point weights make blocks with the same type and live mask
interchangeable, so only their multiplicities need to be stored.
\item Optimize all block interleavings by an exact count-vector DP,
using exact supplied-query codes on positive branches.
\item Charge the clearing prefix against the witness-spine rank
lower bound. Large enough blocks make that charge at most
$\varepsilon$ times the optimum.
\end{enumerate}
\end{proofstrategy}

\subsection{Supporting propositions and the algorithm}

Write $N=|P|=n-1$, $G(T)=\sum_{h\in\HH}d_T(h)$, and
$\OPT_G=\min_TG(T)=n\OPT_E$. We use a positive integer block length
$\ell\ge\max\{2b,b(b-1)\}$ and choose its value in the final
calibration. When $N<\ell$, exact subset Bellman optimization solves the
instance. Assume henceforth that $N\ge\ell$.

The purpose of clearing is to make every surviving query local.
Short gaps suffice because a shallow cap cannot cross a fully
cleared gap and still reach points on both sides.

\begin{proposition}[Clearing gaps isolates small residual blocks]
\label{prop:uniform-cap-prefix}
Partition the convex cyclic order into pieces of length $\ell$ and
mark the first $b$ points of each piece, or the whole piece if shorter.
The number $r$ of marked points satisfies
\begin{equation}\label{eq:uniform-cap-deletion-count}
 r\le b\lceil N/\ell\rceil\le\frac{2bN}{\ell}.
\end{equation}
A legal supplied-query prefix clears all marked points in at most
$r$ spine queries and has total raw contribution
\begin{equation}\label{eq:uniform-cap-prefix-cost}
 G_{\rm prefix}\le r(N+1)+b(b-1)r\le2r(N+1).
\end{equation}
Every remaining cap acts inside one residual block of size at most
$\ell$.
\end{proposition}

The global state records each block's available splits and live
positions. Uniform weights then let us forget a block's identity.

\begin{proposition}[Block types retain all continuation information]
\label{prop:uniform-cap-types}
A residual block's type is its size and its family $\mathcal F_T$ of
nonempty supplied interval masks. There are at most $2^{b\ell}$ types
and at most $d_{b,\ell}$ pairs $(T,R)$ of a type and a nonempty live
mask. Under uniform weights, blocks with the same $(T,R)$ can be
exchanged without changing the optimum continuation cost. Each mask
retains a realizing query for every physical block of its type.
\end{proposition}

Let $c_{T,R}$ count blocks in local state $(T,R)$, and put
$W(c)=1+\sum_{T,R}c_{T,R}|R|$, including the distinguished hypothesis.
For each type, use Algorithm~\ref{alg:exact-tree} to compute the exact
raw local identification cost $g_T(A)$ for every subset $A$ using
its supplied masks; $g_T(A)=0$ for $|A|\le1$. The positive branch
has no distinguished witness. Let $e_{T,R}$ be a unit count vector,
with $e_{T,\varnothing}=0$.

The algorithm must still choose which block to query next. The
count vector preserves these choices, while the local tables handle
the positive exits that no longer interact with other blocks.

\begin{proposition}[The residual count-vector recurrence is exact]
\label{prop:uniform-cap-count-dp}
The optimum residual raw cost is
\begin{equation}\label{eq:uniform-cap-count-dp}
 \begin{aligned}
 F(0)&=0,\\
 F(c)&=W(c)+\min_{\substack{c_{T,R}>0,\ q\in\mathcal F_T\\
                          A=R\cap q\ne\varnothing}}
 \left\{g_T(A)+F\bigl(c-e_{T,R}+e_{T,R\setminus A}\bigr)\right\}.
 \end{aligned}
\end{equation}
The DP has at most $(N+1)^{d_{b,\ell}}$ states, polynomial bit
complexity, and a supplied-query reconstruction realizing its value.
The recurrence includes every adaptive interleaving of the blocks.
\end{proposition}

Clearing buys exact residual optimization at an extra prefix cost.
The witness path already incurs substantial rank cost when caps are
small, so this prefix can be charged against the full optimum.

\begin{proposition}[The clearing charge is small relative to the optimum]
\label{prop:uniform-cap-charge}
The exact residual optimum is at most $\OPT_G$, and
\begin{equation}\label{eq:uniform-cap-rank-lower}
 \OPT_G\ge\sum_{i=1}^N\lceil i/b\rceil
          \ge\frac{N(N+1)}{2b}.
\end{equation}
The prefix followed by the exact residual tree therefore has
$G_{\rm alg}\le(1+8b^2/\ell)\OPT_G$.
\end{proposition}

\begin{algorithm}[H]
\caption{Uniform expected-depth PTAS for bounded convex-position caps}
\label{alg:uniform-cap-ptas}
\KwIn{Uniform hypotheses $P\cup\{p\}$ in convex position; supplied
caps missing $p$, each of size at most fixed $b$; $0<\varepsilon\le1$.}
\KwOut{A supplied-query tree of expected cost at most
$(1+\varepsilon)\OPT_E$.}
$N\gets|P|$; use $\ell$ from
Section~\ref{sec:uniform-cap-calibration}\;
\If{$N<\ell$}{\KwRet $\ExactTree(P\cup\{p\},\mathrm E)$ with raw unit weights\;}
Partition the convex cyclic order of $P$ into pieces of length
$\ell$; mark their first $\min\{b,\text{piece size}\}$ points as $D$\;
$U\gets P$; start an empty miss-spine prefix\;
\While{$D\cap U\ne\varnothing$}{
 Choose $x\in D\cap U$ and a supplied cap $q$ containing $x$\;
 $A\gets U\cap q$\;
 Append $q$ to the prefix; attach a legal local code for $A$
 of height at most $|A|-1$ on its hit side\;
 $U\gets U\setminus A$; continue on the miss side\;
}
Intersect the unmarked pieces with $U$ to obtain the residual blocks\;
Store each physical block's map from its masks to realizing supplied
queries; group identical mask families as types $T$ and compute every
local cost $g_T(A)$ by \ExactTree\;
Compute the initial count vector $c$ of the full residual block masks\;
Evaluate \eqref{eq:uniform-cap-count-dp}, memoizing minimizing transitions\;
Backtrack from $c$: select a physical block in the chosen state,
query that block's stored realizing line, and translate the local
positive code using the same block's query map\;
Attach the reconstructed continuation to the prefix's miss leaf\;
\KwRet the resulting tree\;
\end{algorithm}


\begin{proof}[Proof of Theorem~\ref{thm:uniform-cap-ptas}]
Choose $\ell$ by Section~\ref{sec:uniform-cap-calibration}.
If $N<\ell$, exact subset optimization takes time depending only on
$b,\varepsilon$, apart from polynomial input processing. Otherwise,
Proposition~\ref{prop:uniform-cap-prefix} constructs the legal prefix,
and Propositions~\ref{prop:uniform-cap-types}
and~\ref{prop:uniform-cap-count-dp} give an exact supplied-query
residual tree with the stated time bound.
Proposition~\ref{prop:uniform-cap-charge} and the calibration yield
$G_{\rm alg}\le(1+\varepsilon)\OPT_G$. Divide by $n$ to obtain the
expected-depth ratio. No step introduces an unsupplied singleton test.
\end{proof}

\subsection{Technical proofs of the supporting propositions}

\begin{proof}[Proof of Proposition~\ref{prop:uniform-cap-prefix}]
There are two separate obligations: show that the marked gaps
isolate the surviving blocks, and clear those gaps using only supplied
queries. We begin with the geometric separation property.

\begin{claim}[Marked gaps prevent caps from crossing blocks]
\label{clm:uniform-cap-gap}
Every residual cap meets at most one unmarked block.
\end{claim}
\begin{proof}
Each cap is a cyclic interval in the convex order of $P$: a line
crosses the boundary of its convex polygon at most twice. Partition
that order into consecutive pieces of length $\ell$, with a possibly
shorter last piece. Mark the first $b$ points of each piece, or the
whole piece when its size is at most $b$. Let $D$ be the marked set
and $r=|D|$. There are $\lceil N/\ell\rceil$ pieces and each contributes
at most $b$ marked points. Since $N/\ell\ge1$,
\[
 r\le b\lceil N/\ell\rceil
   \le b(N/\ell+1)\le2bN/\ell,
\]
which proves \eqref{eq:uniform-cap-deletion-count}.
The middle inequality uses $\lceil x\rceil\le x+1$, and the last
uses $N/\ell\ge1$ to absorb the additional piece.
The unmarked portions are blocks of at most $\ell$ points. Between
neighboring nonempty blocks, including across the cyclic wraparound,
there are at least $b$ marked points. A cyclic cap of size at most $b$
cannot meet two such blocks, because it would also contain an entire
intervening gap and hence at least $b+2$ points.

\end{proof}
Supplied caps containing the remaining marked points clear them
legally; removing additional unmarked points can only shrink the
residual blocks.

\begin{claim}[Supplied queries clear the marked set with controlled cost]
No private or arbitrary singleton query is required by the prefix.
\end{claim}
\begin{proof}
While a marked point $x$ remains on the miss spine, query any supplied
cap containing $x$. Such a cap exists because the supplied tests
distinguish $x$ from $p$. Identify its first-hit group by a proper
local supplied-query tree, and continue on the miss side. Separation
gives a local height at most $|A|-1\le b-1$ for each hit group $A$:
repeatedly query a line separating two remaining hypotheses. Each
spine query clears a new marked point, so at most $r$ are used.
Additional unmarked points may leave through these exits. The
residual portions of the blocks still have size at most $\ell$, and
Claim~\ref{clm:uniform-cap-gap} therefore ensures that every residual
cap still lies in one block: clearing extra points cannot create a
new hit across a marked gap.

Each spine node queries at most $N+1$ live hypotheses, so at most
$r$ such nodes contribute $r(N+1)$ to the raw sum of depths.
Its disjoint positive exits contain at most $br$ points in total,
each with at most $b-1$ additional local queries. Therefore
\[
 G_{\rm prefix}\le r(N+1)+br(b-1).
\]
Since $b(b-1)\le\ell\le N$, the second term is at most
$rN\le r(N+1)$. This gives
$G_{\rm prefix}\le2r(N+1)$, as required.

\end{proof}
\end{proof}

\begin{proof}[Proof of Proposition~\ref{prop:uniform-cap-types}]
Relabel the points of each residual block by their positions. Its
type $T$ is its size together with the family $\mathcal F_T$ of all
distinct nonempty supplied-query masks on that block. No cap can
wrap around a block's ends, since it would contain an entire marked
gap. Deleting additional points preserves interval order. Thus every
mask is an ordinary interval of at most $b$ positions. For each possible left endpoint there are at most $b$ possible
interval lengths. Padding to $\ell$ positions therefore gives at
most $b\ell$ masks, hence at most $2^{b\ell}$ subsets of masks
that could form a family. The union of a family's masks is its full block,
since every point is separated from $p$, so the family also determines
the block size. A local
state $(T,R)$ records a nonempty live mask $R$ in a block of type $T$.
For each type there are at most $2^\ell-1$ nonempty live masks,
so the total number of local states is at most
$2^{b\ell}(2^\ell-1)\le2^{(b+1)\ell}=d_{b,\ell}$, using
$2^\ell-1\le2^\ell$.

Two blocks of the same type have identical available splits after
relabelling positions. Store a realizing supplied query separately for
each physical block and mask, so every chosen split uses a line valid
on that block. Each realizing query misses all other residual blocks.
Uniform costs depend only on their live
cardinalities, so exchanging blocks of the same $(T,R)$ preserves
every continuation cost. The count vector therefore retains all needed
information.

\end{proof}

\begin{proof}[Proof of Proposition~\ref{prop:uniform-cap-count-dp}]
Exactness requires both directions: every adaptive tree must obey
the recurrence, and each recurrence choice must have a legal supplied
query realization. Locality supplies the common reason for both.
Any query that is constant on the current live hypotheses can be
deleted: every hypothesis continues to the same child, and removing
the query shortens each of their depths by one. Thus we may suppress
such queries and analyze a nonconstant root query.

\begin{claim}[Every nonconstant adaptive query gives one recurrence transition]
\label{clm:uniform-cap-transition}
A positive branch is local; the miss branch changes only one block state.
\end{claim}
\begin{proof}
By Proposition~\ref{prop:uniform-cap-prefix}, a nonconstant query
hits only one residual block. A query at state $c$ contributes one unit of depth for each live
hypothesis, hence raw cost $W(c)$. Its hit group lies in one block
with local state $(T,R)$ and equals $A=R\cap q\ne\varnothing$.
The miss branch retains $p$ and all other blocks, and changes this
block to $R\setminus A$. The two branch costs are bounded below
by $g_T(A)$ and $F(c-e_{T,R}+e_{T,R\setminus A})$ respectively.
Thus every adaptive tree has cost at least the recurrence value,
by induction on the number of live nonwitness points.

Conversely, for a minimizing transition choose a physical block in
state $(T,R)$, which exists because $c_{T,R}>0$, and its stored
supplied query realizing $q$. The query misses every other residual
block and is proper because its hit group is nonempty and $p$ is
on the miss side. Attach the exact local code of cost $g_T(A)$ and
the inductively reconstructed miss continuation. Their total cost is
$W(c)+g_T(A)+F(c-e_{T,R}+e_{T,R\setminus A})$, attaining the
recurrence. Its live nonwitness count decreases by $|A|\ge1$,
so the induction and the DP are acyclic. The choice of block at every
transition remains unrestricted, covering all adaptive interleavings.
When no live block remains, only $p$ remains and $F(0)=0$.

\end{proof}
Once these transitions are exact, the remaining issue is their
representation. Counting interchangeable block states gives the size
bound, while backtracking retains the physical queries needed for
one output tree.

\begin{claim}[The compressed state space has the stated bit complexity]
Only counts, supplied masks, and integral costs are needed.
\end{claim}
\begin{proof}
There are at most $d_{b,\ell}$ count coordinates, each between
zero and $N$, so the enclosing count cube has at most
$(N+1)^{d_{b,\ell}}$ states. For a coordinate with positive count,
there are at most $b\ell$ supplied masks to try. Each local Bellman
table concerns at most $2^\ell$ subsets; these bounds depend only
on $b,\ell$. Constructing the type and physical-query maps requires
polynomial input processing.

Every reachable state contains at most $n$ actual hypotheses. A proper
identifying tree for it has height at most $n-1$, giving raw cost
at most $n(n-1)$. All weights are one, so the DP costs are integers
with $O(\log n)$ bits on reachable states. Storing a minimizing
transition and a physical realizing query at each used state suffices
for backtracking; no list of all adaptive trees is stored.

\end{proof}
\end{proof}

\begin{proof}[Proof of Proposition~\ref{prop:uniform-cap-charge}]
Restrict an optimal full tree to the residual hypotheses, suppressing
constant-answer nodes. Lemma~\ref{lem:restriction} shows that this
can only decrease depths; Proposition~\ref{prop:uniform-cap-count-dp}
therefore makes the exact residual cost at most
$\OPT_G$. For the lower bound, order the nonwitness hypotheses by their
first-hit position on the witness spine. The first $j$ queries can
remove at most $bj$ of them, so the $i$th hypothesis in this order
has first-hit position at least $\lceil i/b\rceil$ and leaf depth
at least that large. Ignoring the nonnegative witness and continuation
costs yields
\[
 \OPT_G\ge\sum_{i=1}^N\lceil i/b\rceil
 \ge\sum_{i=1}^N i/b=\frac{N(N+1)}{2b}.
\]
The second inequality uses $\lceil i/b\rceil\ge i/b$ term by term,
and the equality sums the first $N$ integers.
Write $G_{\rm res}$ for the exact residual cost. The restriction
argument gives $G_{\rm res}\le\OPT_G$, while the prefix estimate
gives $G_{\rm alg}\le2r(N+1)+G_{\rm res}$. Consequently
\begin{align*}
 \frac{G_{\rm alg}-\OPT_G}{\OPT_G}
 &\byref{Prop.}{prop:uniform-cap-prefix}{\le}\frac{2r(N+1)}{\OPT_G}\\
 &\le\frac{2r(N+1)}{N(N+1)/(2b)}
       =\frac{4br}{N}\\
 &\byref{Eq.}{eq:uniform-cap-deletion-count}{\le}\frac{4b}{N}\frac{2bN}{\ell}
       =\frac{8b^2}{\ell}.
\end{align*}
The first inequality uses residual optimality, the second the rank
lower bound, and the last the marked-point count.
Dividing by $n$ proves the expected-depth guarantee.
\end{proof}

\subsection{Calibration of the block length}
\label{sec:uniform-cap-calibration}

\begin{parameterchoice}
For $0<\varepsilon\le1$, choose
\[
 \ell=\max\left\{2b,\ b(b-1),\
                 \left\lceil\frac{8b^2}{\varepsilon}\right\rceil\right\}.
\]
The first two bounds ensure gap separation and the prefix estimate;
the last gives $8b^2/\ell\le\varepsilon$. For fixed $b,\varepsilon$,
$\ell$ and $d_{b,\ell}$ are constants, yielding the claimed PTAS time.
\end{parameterchoice}


The proof uses only the cyclic-interval property of the caps. In
particular it also applies when the nonwitness points alone are in
convex position. For fixed $b$, a constant-factor hardness reduction
through uniform one-layer shallow-cap instances would contradict
this PTAS unless $\mathrm P=\mathrm{NP}$. A construction based on
bounded-size caps must therefore prove interactions beyond this
cyclic-interval subclass to establish such a gap.
